\section{Related Work}
\label{sec:related}

\subsection{Machine Learning and Meta-Labeling in Trading}

Gu, Kelly, and Xiu \cite{gu2020empirical} show that tree ensembles and neural networks improve return forecasts and
the portfolios built from them, largely by capturing nonlinear interactions among momentum, liquidity, and volatility
signals. Their object is the forecast, not the decision it feeds. Meta-labeling \cite{lopezdeprado2018afml} addresses
the decision directly: a primary model chooses the side of a trade, and a secondary classifier, trained on realized
outcomes, decides whether to take it and how large to make it. Joubert \cite{joubert2022metalabeling} formalizes
meta-labeling as a layer that filters false-positive signals and sizes positions on top of a primary strategy. Our
meta-labeling and position-sizing components implement exactly these two uses, with a fixed rule rather than a learned
model as the primary; in the taxonomy of Section~\ref{sec:taxonomy} both are \emph{narrowing} components, because
they can only reject or re-weight decisions the rule has already made.

\subsection{Hybrid Systems and Module Ablations}

Hybrid systems that combine technical rules, ML, and other signals typically report aggregate improvements over a
rule-only baseline without decomposing which sub-component is responsible; a recent regime-adaptive equity system that
integrates technical analysis, ML, and financial sentiment is a representative example
\cite{generatingalpha2026hybrid}. The closest methodological precedent we found is a deep-learning multi-day-turnover
pipeline for Chinese A-shares \cite{turnover2025ablation}, which ablates its modules one at a time and reports each
module's marginal contribution to annualized return. MBCA shares that spirit but differs in three ways: the
components are insertion points of one fixed rule rather than modules of a proprietary architecture, every component
receives the same feature and model budget, and every component is tested out of sample at a pre-specified freeze.
Volatility-managed portfolios \cite{moreira2017volmanaged} scale exposure by realized volatility at the portfolio
level; our sizing component instead scales each trade by a predicted probability of success.

\subsection{Technical Trading Rules and Data Snooping}

Moving-average and breakout rules have a long empirical record \cite{brock1992simple,park2007technical}, and
time-series momentum is documented across asset classes and more than a century of data
\cite{moskowitz2012momentum,hurst2017century}. Their apparent profitability, however, is highly sensitive to how many
rules were searched. White's reality check \cite{white2000reality} and its application to technical rules by
Sullivan, Timmermann, and White \cite{sullivan1999data} remove much of it once the full rule universe is accounted
for; Hansen's test of superior predictive ability \cite{hansen2005spa} sharpens that comparison, and Harvey, Liu, and
Zhu \cite{harvey2016cross} argue for significance thresholds far above conventional ones. Study~2 therefore evaluates
canonical rules chosen and parameterized before any data were seen, rather than searched ones.

\subsection{Inference on Sharpe Ratios and Backtest Overfitting}

Sharpe ratios are estimated with substantial error, and serial correlation distorts naive annualization
\cite{lo2002sharpe}. Ledoit and Wolf \cite{ledoit2008robust} propose a studentized test of the difference between
two Sharpe ratios based on a circular block bootstrap \cite{politis1992circular}, which respects both the pairing of
the two return series and their serial dependence. Selection across many backtests inflates the best observed Sharpe
ratio; the deflated Sharpe ratio \cite{bailey2014deflated} and the probability of backtest overfitting
\cite{bailey2017pbo} quantify that bias. Across many tests, the false-discovery rate \cite{benjamini1995fdr} and Holm's
family-wise procedure \cite{holm1979simple} control multiplicity, and pre-registration \cite{nosek2018preregistration}
separates confirmatory tests from exploratory ones.

\subsection{Positioning}

We found no work that holds a rule-based strategy, its data, features, model class, and evaluation protocol fixed
while sweeping which single decision a learned model replaces, that names this constraint as a reusable
methodology, and that reports every resulting configuration rather than the best one. MBCA is complementary to
data-snooping corrections: those account for the search \emph{across} rules, whereas MBCA controls what varies
\emph{within} one rule. This paper is the first full application of the methodology.
